\section{Proof of the exact ordinary-gap range}
\label{gapregime:proof}

\paragraph{A scalar linear extremal problem.}
Write $K_g(t)=H_0+H_1gt$ on $[0,1]$.  The strictly convex function
\[
 D_g(a)=\int_0^1|a-t|K_g(t)\,dt-ma,\qquad 0\le a\le1,
\]
satisfies
\[
 D_g'(a)=2H_0a+H_1ga^2-H_0-H_1g/2-m,
 \qquad D_g''(a)=2K_g(a)>0.
\]
Its unique minimizer is $a=\tau_g$ from
\eqref{gapregime:switch}.  At $g=0,H_0=m$ this is the endpoint
$\tau_0=1$; otherwise it is interior.  At $a=1/2$ the derivative
is negative, so $1/2<\tau_g\le1$.
Direct integration, followed by the switch equation, gives
\begin{align*}
 \min_{0\le a\le1}D_g(a)
 &=H_0(\tau_g^2-\tau_g+1/2)
   +H_1g(\tau_g^3/3-\tau_g/2+1/3)-m\tau_g\\
 &=\mathcal S(g).
\end{align*}
If a measurable acceleration $a(t)$ satisfies
$|a(t)|\le K_g(t)$ and $\int_0^1a(t)\,dt=m$, then
\begin{equation}\label{gapregime:dual-bound}
 \int_0^1(1-t)a(t)\,dt-m
 =\int_0^1(\tau_g-t)a(t)\,dt-m\tau_g
 \le\mathcal S(g).
\end{equation}
Equality is obtained by $a(t)=K_g(t)$ before $\tau_g$ and
$a(t)=-K_g(t)$ afterwards.  The switch equation is precisely its
integral constraint.  This is an elementary extremal problem over
bounded functions, not an additional optimization oracle.

\paragraph{First variation at zero radius.}
Fix $g>0$ and $H_1>0$.  Let $u_{\rho,\tau}$ solve the normalized
two-arc system
\[
 u''=\rho\sigma_\tau(t)(H_0+H_1u),\qquad
 u(0)=0,\quad u'(0)=g-m\rho,
\]
where $\sigma_\tau$ equals $+1$ before $\tau$ and $-1$ afterwards.
At $\rho=0$ its solution is $u=gt$, independently of $\tau$.
For $\rho\ne0$, the terminal slope equation $u'(1)=g$ is equivalent
to
\[
 R(\rho,\tau):=-m+
       \int_0^1\sigma_\tau(t)(H_0+H_1u_{\rho,\tau}(t))\,dt=0.
\]
The explicit two-arc solutions make $R$ smooth near
$(0,\tau_g)$, and
\[
 R(0,\tau_g)=0,\qquad
 \partial_\tau R(0,\tau_g)=2K_g(\tau_g)>0.
\]
Because $g>0,H_1>0$, this switch is interior.  The implicit function
theorem supplies a smooth branch $\tau(\rho)$ from $\tau_g$.
For small positive $\rho$, uniqueness of the fixed-radius switch
identifies its terminal loss with $V(g,\rho)$.
Twice integrating the equation gives
\begin{equation}\label{gapregime:variation}
 \begin{aligned}
 V(g,\rho)
 &=g+\rho\left\{-m+
       \int_0^1(1-t)\sigma_{\tau_g}(t)K_g(t)\,dt\right\}
         +O(\rho^2)\\
 &=g+\rho\mathcal S(g)+O(\rho^2).
 \end{aligned}
\end{equation}
In particular, $\mathcal S(g)>0$ implies
$\mathcal B(g)>g$.

\paragraph{Excluding a crossing when $\mathcal S(g)<0$.}
Equation~\eqref{gapregime:variation} gives $V(g,\rho)<g$ for all
sufficiently small positive $\rho$.  Suppose that this inequality
fails somewhere on $(0,g/m]$.  By continuity there is a first
positive $\rho_*$ with $V(g,\rho_*)=g$.
Let $u$ be its normalized extremal two-arc trajectory.  Its derivative
starts at $g-m\rho_*<g$, increases on the first arc, and decreases
on the second arc to $g$.  The switch is interior and the peak
derivative exceeds $g$.  Hence $u'(t)-g$ changes sign exactly once,
from negative to positive, before ending at zero.
Since $u(0)=0$ and $u(1)=g$, it follows that
\begin{equation}\label{gapregime:below-line}
 u(t)<gt\qquad(0<t<1).
\end{equation}
In particular $|u''(t)|/\rho_*\le K_g(t)$ almost everywhere, while
$\int_0^1u''(t)/\rho_*\,dt=m$.  Apply
\eqref{gapregime:dual-bound} with $a=u''/\rho_*$ to obtain
\[
 0=u(1)-g\le\rho_*\mathcal S(g)<0,
\]
a contradiction.  Thus $V(g,\rho)<g$ for every positive radius,
and the definition $V(g,0)=g$ proves $\mathcal B(g)=g$.

\paragraph{Monotonicity of the first-variation coefficient.}
The formula \eqref{gapregime:tau} is smooth for $g\ge0$.
Differentiating $D_g(\tau_g)$ eliminates the switch derivative,
because $D_g'(\tau_g)=0$, including at $g=0,H_0=m$.
Therefore
\[
 \mathcal S'(g)=H_1\int_0^1|\tau_g-t|t\,dt
 =H_1(\tau_g^3/3-\tau_g/2+1/3)>0.
\]
For $a\in[1/2,1]$, the function
$j(a)=a^3/3-a/2+1/3$ has its unique minimum at $a=1/\sqrt2$.
Its maximum on this interval is $j(1)=1/6$.  Thus
\begin{equation}\label{gapregime:derivative-bounds}
 \frac{H_1}{3}(1-1/\sqrt2)
       \le\mathcal S'(g)\le\frac{H_1}{6}.
\end{equation}
At $g=0$, substitution of
$\tau_0=(H_0+m)/(2H_0)$ proves
\eqref{gapregime:initial}.  The strictly positive lower derivative
bound shows that $\mathcal S(g)\to\infty$.
These facts prove existence and uniqueness of a positive zero
exactly when $m\le H_0<(1+\sqrt2)m$, and integrating
\eqref{gapregime:derivative-bounds} proves
\eqref{gapregime:threshold-bracket}.

\paragraph{The zero coefficient and strictness.}
If $g>0$ and $\mathcal S(g)=0$, strict monotonicity gives
$\mathcal S(g_n)<0$ for any sequence $0<g_n\uparrow g$.
The preceding noncrossing argument gives
$\mathcal B(g_n)=g_n$.  Continuity of $\mathcal B$ from
Theorem~\ref{boundarygap:main} yields $\mathcal B(g)=g$.
This argument concerns a positive zero only; the case $g=0$ is
already part of the envelope definition.

If some positive radius attained $V(g,\rho)=g$, its trajectory
would satisfy \eqref{gapregime:below-line}.  Since $H_1>0$,
\[
 |u''(t)|/\rho=H_0+H_1u(t)<K_g(t)
 \quad\hbox{for almost every }t\in(0,1).
\]
The integral in \eqref{gapregime:dual-bound} is then strictly
smaller than $\mathcal S(g)$: the multiplier
$|\tau_g-t|$ is positive except at one point.  This gives
$0=V-g<\rho\mathcal S(g)=0$, a contradiction.
Together with the negative-coefficient case, this proves strictness
at every positive radius whenever $\mathcal S(g)\le0$.

\paragraph{Radius monotonicity.}
Take $0<\rho_1<\rho_2\le g/m$ and put $\lambda=\rho_1/\rho_2$.
Let $u_2$ be the normalized extremal trajectory for $\rho_2$.
Its derivative increases and then decreases to $g$, while
$u_2'(0)=g-m\rho_2<g$ and $u_2(1)=V(g,\rho_2)<g$.
Thus $u_2(t)-gt$ first decreases and then increases to a negative
endpoint; in particular $u_2(t)\le gt$ on $[0,1]$.
Define
\[
 v(t)=\lambda u_2(t)+(1-\lambda)gt.
\]
Then $v(0)=0$, $v'(0)=g-m\rho_1$, and $v'(1)=g$.
The normalized shifted-RSI inequality is preserved:
\[
 v'(t)\ge g-m\rho_1+m\rho_1t.
\]
Moreover $v\ge u_2$ and, almost everywhere,
\[
 |v''|\le\lambda\rho_2(H_0+H_1u_2)
          \le\rho_1(H_0+H_1v).
\]
The fixed-radius energy comparison in Appendix~\ref{boundarygap:proof}
therefore gives
\begin{equation}\label{gapregime:mixture}
 V(g,\rho_1)\ge v(1)
       =\lambda V(g,\rho_2)+(1-\lambda)g>V(g,\rho_2).
\end{equation}
The comparison applies to this $C^1$, piecewise $C^2$ trajectory:
its derivative is absolutely continuous, so both energy integrations
and the change of variable in that proof remain valid with the
almost-everywhere curvature bound.  No smooth extremizer is needed.
Rearranging \eqref{gapregime:mixture} proves
\eqref{gapregime:deficit-ratio}.  Letting $\rho_1\downarrow0$
and using \eqref{gapregime:variation} gives
\eqref{gapregime:tangent-upper}.  The value $V(g,0)=g$ and the
strict inequality at every positive radius include zero in the
strict monotonicity statement.

\paragraph{The exact threshold when $H_0=m$.}
Take
\[
 g_{\rm c}=(2+\sqrt6)m/H_1,
 \qquad \tau_{g_{\rm c}}=2-\sqrt6/2.
\]
Direct substitution gives both
\[
 (2+\sqrt6)\tau_{g_{\rm c}}^2+2\tau_{g_{\rm c}}
      =2+(2+\sqrt6)/2
\]
and
\[
 \frac12-\tau_{g_{\rm c}}^2
   +\frac{2+\sqrt6}{3}(1-2\tau_{g_{\rm c}}^3)=0.
\]
The number $\tau_{g_{\rm c}}$ lies in $(1/2,1)$, so it is the
unique switch from \eqref{gapregime:switch}.  Strict monotonicity of
$\mathcal S$ proves that this is its unique positive zero.
The special case $H_1=0$ does not require this argument and follows
directly from \eqref{boundarygap:constant}.
